% Official solution ("Problem 3 A. Marking scheme -- Thermodynamic test, 25 points",
% 2014_da_sol.pdf pp.29-35). Decimal commas kept as printed.

\begin{description}
\item[a)]
The pressure and the density of the atmosphere increase with the decrease of the altitude.
The pressure at the surface of the planet is $p_{\max}$.

In figure 1 is represented a thin atmospheric layer with thickness $\Delta h$ near the surface
of the planet P$_1$ where the atmosphere density can be considered $\rho_{\max}$. Inside this
layer the pressure and the temperature are variable, reaching the values $p_{\max}$ and $T_0$
near the surface.
% FIGURE NEEDED: Fig. 1 (2014_da_sol.pdf p.29): thin near-surface layer of thickness
% \Delta h with density \rho_max, probe descending with velocity v, gravity g_0; above it
% a layer of density \rho with values (T; p); labels \Delta p = \rho_0 g_0 \Delta h and
% T_0; p_max = p + \Delta p.

The following relations can be written:
\begin{align*}
p_{\max} &= p + \Delta p;\\
\Delta p &= \rho_{\max}g_0\Delta h,
\end{align*}
$\Delta p$ -- the variation of the pressure due to the altitude variation $\Delta h$, $g_0$
the gravitational acceleration near the surface of planet P$_1$;
\[
\Delta h = \frac{\Delta p}{\rho_{\max}g_0}.
\]
The superior layer where the density is constant $\rho$ is represented too. At its base the
pressure and temperature hav the values $p$ and respectively $T$. % sic: "hav"

For a gas volume inside the near surface layer of the atmosphere, near the surface of the
planet for $\nu$ mol of de gas according with Mendeleev--Klapeyron law % sic: "of de gas"
\begin{align*}
p_{\max}V &= \nu RT_0 = \frac{m}{\mu}RT_0;\\
p_{\max}\mu &= \frac{m}{V}RT_0 = \rho_{\max}RT_0;\\
\rho_{\max} &= \frac{p_{\max}\mu}{RT_0}.
\end{align*}

Just before touching the surface of the planet P$_1$ the radio-probe moves uniformly, so the
distance $\Delta h$ is spend with speed $\vec{\mathrm{v}}$, in time $\Delta t$: % sic: "is spend"
\begin{align*}
\Delta h &= \mathrm{v}\,\Delta t;\\
\Delta h &= \frac{\Delta p}{\rho_{\max}g_0};\quad
\rho_{\max} = \frac{p_{\max}\mu}{RT_0};\\
\frac{\Delta p}{\rho_{\max}g_0} &= \mathrm{v}\,\Delta t;\quad
\mathrm{v} = \frac{\Delta p}{\rho_{\max}g_0\Delta t};\quad
\mathrm{v} = \frac{\Delta p}{\dfrac{p_{\max}\mu}{RT_0}g_0\Delta t};\\
\mathrm{v} &= \frac{RT_0\Delta p}{g_0\mu p_{\max}\Delta t},\quad(1)
\end{align*}
Where accordingly to fig.~2 in the point F $(\tau;\,p_{\max})$, the trigonometric tangent of
$\alpha$ is:
\[
\tan\alpha = \frac{\Delta p}{\Delta t}.
\]
% FIGURE NEEDED: Fig. 2 (2014_da_sol.pdf p.30): the p(u.c.) vs t(s) curve of fig. 3.1 with
% the tangent at the last point F marked, angle \alpha, \tau = 3750 s, p_max and the
% increments \Delta p, \Delta t indicated.

The rapport: % sic: "rapport"
\[
\frac{\Delta p}{p_{\max}\Delta t},
\]
Can be determined in the graph in the point F:
\[
\mathrm{v} = \frac{\Delta p}{p_{\max}\Delta t}\,\frac{RT_0}{g_0\mu};
\]
\[
\frac{\Delta p}{p_{\max}\Delta t}
= \frac{15\,\mathrm{u.c.}}{55\,\mathrm{u.c.}\,250\,\mathrm{s}}
= \frac{3}{2750}\,\frac{1}{\mathrm{s}}
\qquad
\frac{45\,\mathrm{u.c.}}{55\,\mathrm{u.c.}\,750\,\mathrm{s}}
= \frac{3}{2750}\,\frac{1}{\mathrm{s}};
\]
\begin{align*}
\mathrm{v} &= \frac{3}{2750}\,\frac{1}{\mathrm{s}}\,
\frac{8{,}3\,\dfrac{\mathrm{J}}{\mathrm{mol\,K}}\;700\,\mathrm{K}}
     {10\,\dfrac{\mathrm{m}}{\mathrm{s}^2}\;44\,\dfrac{10^{-3}\,\mathrm{kg}}{\mathrm{mol}}}
= 14{,}4\,\frac{\mathrm{m}}{\mathrm{s}};\\
h_0 &= \mathrm{v}\tau;\quad \tau = 3.750\ \mathrm{s};\\
h_0 &= 54.000\,\mathrm{m} = 54\,\mathrm{km}.
\end{align*}
% Dots in 3.750 s and 54.000 m are thousands separators, as in the original.

\item[b)]
Descending from the initial altitude $h_0 = 54\,\mathrm{km}$, with speed
$\mathrm{v} = 14{,}4\,\frac{\mathrm{m}}{\mathrm{s}}$, the radio-probe, after 1000 s arrives
at the altitude $h = 39{,}6\,\mathrm{km}$
\[
t = \frac{h_0 - h}{\mathrm{v}} = \frac{14400\,\mathrm{m}}{14{,}4\,\frac{\mathrm{m}}{\mathrm{s}}}
= 1000\,\mathrm{s}.
\]
On Graph at altitude $h = 39{,}6\,\mathrm{km}$, the atmospheric pressure is
$p_h = 4\,\mathrm{u.c.}$, -- see point C on grah. % sic: "grah"
% FIGURE NEEDED: Fig. 3 (2014_da_sol.pdf pp.31-32): the same p(u.c.) vs t(s) curve with
% point C at p_h and the tangent at C making angle \beta with the t axis.

In this point
\begin{align*}
\tan\beta &= \left(\frac{\Delta p}{\Delta t}\right)_{\!C};\\
\left(\frac{\Delta p}{p\Delta t}\right)_{\!C}
&= \frac{1}{p_h}\left(\frac{\Delta p}{\Delta t}\right)_{\!C}
= \frac{1}{4\,\mathrm{u.c.}}\,\frac{3\,\mathrm{u.c.}}{500\,\mathrm{s}}
= \frac{3}{2000}\,\frac{1}{\mathrm{s}};\\
\left(\frac{p\Delta t}{\Delta p}\right)_{\!C}
&= p_h\left(\frac{\Delta t}{\Delta p}\right)_{\!C}
= 4\,\mathrm{u.c.}\,\frac{500\,\mathrm{s}}{3\,\mathrm{u.c.}}
= \frac{2000}{3}\,\mathrm{s};\\
\mathrm{v} &= \left(\frac{\Delta p}{p\Delta t}\right)_{\!C}\frac{RT_h}{g_0\mu};\\
\mathrm{v} &= \frac{1}{p_h}\left(\frac{\Delta p}{\Delta t}\right)_{\!C}\frac{RT_h}{g_0\mu};\\
T_h &= \frac{\mathrm{v}\,g_0\mu}
{\dfrac{1}{p_h}\left(\dfrac{\Delta p}{\Delta t}\right)_{\!C}R};\quad
T_h = p_h\left(\frac{\Delta t}{\Delta p}\right)_{\!C}\frac{\mathrm{v}\,g_0\mu}{R};\\
T_h &= \frac{2000}{3}\,\mathrm{s}\,
\frac{14{,}4\,\dfrac{\mathrm{m}}{\mathrm{s}}\;10\,\dfrac{\mathrm{m}}{\mathrm{s}^2}\;
      44\,\dfrac{10^{-3}\,\mathrm{kg}}{\mathrm{mol}}}
     {8{,}3\,\dfrac{\mathrm{J}}{\mathrm{mol\,K}}};\\
T_h &= 508{,}9\,\mathrm{K}.
\end{align*}

\item[c)]
By using the same algorithm, corresponding to the points F, D, C, B \c{s}i A in graph on
figure 4: % sic: Romanian "şi" (and) left in the original
% FIGURE NEEDED: Fig. 4 (2014_da_sol.pdf p.32): the linear p(u.c.) vs t(s) graph of the
% planet P2 with the data points A, B, C, D, F marked and the tangent angle \alpha,
% increments \Delta t, \Delta p and \tau at the last point.
\begin{align*}
\mathrm{v} &= \frac{\Delta p}{p_{\max}\Delta t}\,\frac{RT_0}{g_0\mu};\\
\mathrm{v} &= \frac{10\,\mathrm{u.c.}}{50\,\mathrm{u.c.}\,500\,\mathrm{s}}\,
\frac{8{,}3\,\dfrac{\mathrm{J}}{\mathrm{mol\,K}}\;750\,\mathrm{K}}
     {8\,\dfrac{\mathrm{m}}{\mathrm{s}^2}\;44\,\dfrac{10^{-3}\,\mathrm{kg}}{\mathrm{mol}}}
\approx 7\,\frac{\mathrm{m}}{\mathrm{s}};\\
h_0 &= \mathrm{v}\tau = 7\,\frac{\mathrm{m}}{\mathrm{s}}\,2000\,\mathrm{s}
= 14.000\,\mathrm{m} = 14\,\mathrm{km};\\
h_F &= 0;
\end{align*}
\[
\mathrm{F}(p = 50\,\mathrm{u.c.};\ h = 0\,\mathrm{km};\ T = 750\,\mathrm{K});
\]
\begin{align*}
T_D &= p_D\left(\frac{\Delta t}{\Delta p}\right)_{\!D}\frac{\mathrm{v}\,g_0\mu}{R}
= 593{,}7\,\mathrm{K};\\
h_D &= h_0 - 7\,\frac{\mathrm{m}}{\mathrm{s}}\,1500\,\mathrm{s}
= 14\,\mathrm{km} - 10.500\,\mathrm{m} = 3{,}5\,\mathrm{km};
\end{align*}
\[
\mathrm{D}(p = 40\,\mathrm{u.c.};\ h = 3{,}5\,\mathrm{km};\ T = 593{,}7\,\mathrm{K});
\]
\begin{align*}
T_C &= p_C\left(\frac{\Delta t}{\Delta p}\right)_{\!C}\frac{\mathrm{v}\,g_0\mu}{R}
= 445{,}3\,\mathrm{K};\\
h_C &= h_0 - 7\,\frac{\mathrm{m}}{\mathrm{s}}\,1000\,\mathrm{s}
= 14\,\mathrm{km} - 7.000\,\mathrm{m} = 7\,\mathrm{km};
\end{align*}
\[
\mathrm{C}(p = 30\,\mathrm{u.c.};\ h = 7\,\mathrm{km};\ T = 445{,}3\,\mathrm{K});
\]
\begin{align*}
T_B &= p_B\left(\frac{\Delta t}{\Delta p}\right)_{\!B}\frac{\mathrm{v}\,g_0\mu}{R}
= 296{,}8\,\mathrm{K};\\
h_B &= h_0 - 7\,\frac{\mathrm{m}}{\mathrm{s}}\,500\,\mathrm{s}
= 14\,\mathrm{km} - 3.500\,\mathrm{m} = 10{,}5\,\mathrm{km};
\end{align*}
\[
\mathrm{B}(p = 20\,\mathrm{u.c.};\ h = 10{,}5\,\mathrm{km};\ T = 296{,}8\,\mathrm{K});
\]
\begin{align*}
T_A &= p_A\left(\frac{\Delta t}{\Delta p}\right)_{\!A}\frac{\mathrm{v}\,g_0\mu}{R}
= 148{,}4\,\mathrm{K};\\
h_A &= 14\,\mathrm{km};
\end{align*}
\[
\mathrm{A}(p = 10\,\mathrm{u.c.};\ h = 14\,\mathrm{km};\ T = 148{,}4\,\mathrm{K}).
\]

The data are in table below which allows to plot de graph $p = f(h)$ and $T = f(h)$ for
planet P$_2$ figure 5 and figure 6. % sic: "plot de graph"
\begin{center}
\begin{tabular}{|c|c|c|c|c|c|}
\hline
$t$ & 0 & 500 s & 1000 s & 1500 s & 2000 s\\
\hline
$p$ & 10 u.c. & 20 u.c. & 30 u.c. & 40 u.c. & 50 u.c.\\
\hline
$h$ & 14 km & 10,5 km & 7 km & 3,5 km & 0\\
\hline
$T$ & 148,4 K & 296,8 K & 445,3 K & 593,3 K & 750 K\\
\hline
\end{tabular}
\end{center}
% sic: the table prints 593,3 K although the computation above gives 593,7 K.
% FIGURE NEEDED: Fig. 5 (2014_da_sol.pdf p.34): p(u.c.) vs h(km) straight line through the
% points A(14; 10) ... F(0; 50) with labels A-F.
% FIGURE NEEDED: Fig. 6 (2014_da_sol.pdf pp.34-35): T(K) vs h(km) straight line through the
% points A(14 km; 148,4 K) ... F(0; 750 K) with labels A-F.
\end{description}
